% ============================================================================
% preamble.tex — 包、字体、排版、代码环境、自定义彩框
% 《多邻国 Android 面试通关手册》
% ============================================================================

% --- 中文支持 (XeLaTeX) ---
\usepackage{ctex}
\usepackage{fontspec}

% 等宽字体：DejaVu Sans Mono（渲染代码 / 框图制表符；需安装到系统字体）
\setmonofont{DejaVu Sans Mono}[Scale=0.92]

% --- 页面布局 ---
\usepackage[
  a4paper,
  top=2.4cm, bottom=2.4cm,
  left=2.3cm, right=2.3cm
]{geometry}

% --- 数学 ---
\usepackage{amsmath, amssymb, amsthm}
\usepackage{mathtools}

% --- 图表 ---
\usepackage{graphicx}
\usepackage{float}
\usepackage{booktabs}
\usepackage{longtable}
\usepackage{tabularx}
\usepackage{array}
\usepackage{multirow}

% --- 绘图 ---
\usepackage{tikz}
\usetikzlibrary{positioning, calc, arrows.meta, fit, backgrounds, shapes.geometric}

% --- 颜色 ---
\usepackage{xcolor}
\definecolor{duogreen}{HTML}{58CC02}  % Duolingo 品牌绿
\definecolor{duoblue}{HTML}{1CB0F6}   % Duolingo 辅助蓝
\definecolor{duoorange}{HTML}{FF9600} % Duolingo streak 橙
\definecolor{duored}{HTML}{FF4B4B}    % Duolingo hearts 红
\definecolor{abred}{HTML}{58CC02}     % 兼容别名 -> 品牌主色
\definecolor{abdark}{HTML}{4B4B4B}
\definecolor{abgray}{HTML}{777777}
\definecolor{codebg}{HTML}{F7F7F9}
\definecolor{codekw}{HTML}{0033B3}
\definecolor{codestr}{HTML}{067D17}
\definecolor{codecmt}{HTML}{8C8C8C}
\definecolor{codetype}{HTML}{8000A0}

% --- 超链接与交叉引用 ---
\usepackage{hyperref}
\hypersetup{
  colorlinks=true,
  linkcolor=abred!70!black,
  citecolor=green!40!black,
  urlcolor=cyan!55!black,
  bookmarksnumbered=true,
}
\usepackage[nameinlink]{cleveref}

% --- 代码高亮 (listings, 纯 LaTeX, 无需 python) ---
\usepackage{xstring}
\usepackage{listings}
\lstset{
  basicstyle=\ttfamily\small,
  backgroundcolor=\color{codebg},
  keywordstyle=\color{codekw}\bfseries,
  stringstyle=\color{codestr},
  commentstyle=\color{codecmt}\itshape,
  numberstyle=\tiny\color{abgray},
  numbers=left, numbersep=8pt,
  frame=single, framesep=5pt, rulecolor=\color{gray!30},
  breaklines=true, breakatwhitespace=false,
  showstringspaces=false,
  tabsize=2,
  columns=fullflexible,
  keepspaces=true,
  xleftmargin=14pt, xrightmargin=4pt,
  aboveskip=8pt, belowskip=8pt,
}
% Kotlin 语言定义（listings 内置 Kotlin 支持不完整，这里自定义）
\lstdefinelanguage{Kotlin}{
  keywords={
    package, import, class, interface, object, data, sealed, enum, annotation,
    fun, val, var, const, lateinit, by, lazy, init, constructor, companion,
    override, open, abstract, final, private, protected, public, internal,
    suspend, inline, noinline, crossinline, reified, operator, infix, tailrec,
    vararg, out, in, where, typealias, actual, expect,
    if, else, when, for, while, do, return, break, continue, throw, try, catch,
    finally, is, as, this, super, null, true, false, it, typeof
  },
  morekeywords=[2]{
    launch, async, await, runBlocking, withContext, coroutineScope,
    supervisorScope, delay, flow, collect, collectLatest, emit, cancel,
    ensureActive, yield, remember, mutableStateOf, LaunchedEffect,
    DisposableEffect, SideEffect, rememberCoroutineScope,
    collectAsState, collectAsStateWithLifecycle, repeatOnLifecycle,
    stateIn, shareIn, flowOn, distinctUntilChanged, debounce, flatMapLatest,
    combine, map, filter, forEach, let, apply, run, also, takeIf, require
  },
  morekeywords=[3]{
    Int, Long, Double, Float, Boolean, String, Char, Unit, Any, Nothing,
    List, MutableList, Map, MutableMap, Set, MutableSet, Array, Result,
    Flow, StateFlow, MutableStateFlow, SharedFlow, MutableSharedFlow, Channel,
    LiveData, MutableLiveData, ViewModel, Job, Deferred, CoroutineScope,
    Dispatchers, GlobalScope, viewModelScope, lifecycleScope, WorkManager,
    Context, Activity, Fragment, RecyclerView, ViewHolder, Composable
  },
  sensitive=true,
  morecomment=[l]{//},
  morecomment=[s]{/*}{*/},
  morestring=[b]",
  morestring=[b]',
}

\lstdefinestyle{kt}{
  language=Kotlin,
  keywordstyle=\color{codekw}\bfseries,
  keywordstyle=[2]\color{duoblue!75!black},
  keywordstyle=[3]\color{codetype},
}
\lstnewenvironment{kotlin}{\lstset{style=kt}}{}
% 带 bug 的代码（红色边框，视觉上与正确代码区分）
\lstnewenvironment{ktbug}{%
  \lstset{style=kt, rulecolor=\color{duored!60}, backgroundcolor=\color{duored!4}}%
}{}
% 修正后的代码（绿色边框）
\lstnewenvironment{ktfix}{%
  \lstset{style=kt, rulecolor=\color{duogreen!70!black}, backgroundcolor=\color{duogreen!5}}%
}{}

% C++ 别名（保留，个别对照场景使用）
\lstdefinestyle{cpp}{
  language=C++,
  morekeywords={auto, vector, string, unordered_map, unordered_set, map, set,
    priority_queue, queue, deque, pair, size_t, int64_t, nullptr, override,
    constexpr, string_view},
  emph={int, long, bool, char, double, void, struct, class},
  emphstyle=\color{codetype},
}
\lstnewenvironment{cpp}{\lstset{style=cpp}}{}

% --- 彩色定理框 ---
\usepackage[most]{tcolorbox}
\tcbuselibrary{breakable, skins}

% ============================================================================
% 自定义结构化彩框：本书面经的核心“设计思想”载体
% ============================================================================

% 设计思想 / 核心心法（红框）
\newtcolorbox{idea}[1][核心思想]{
  enhanced, breakable, arc=3pt,
  before skip=10pt, after skip=10pt,
  colback=abred!5, colframe=abred!80!black,
  boxrule=0.8pt, left=8pt, right=8pt, top=6pt, bottom=6pt,
  fonttitle=\bfseries, title={【设计思想】#1},
}

% 陷阱 / 易错点（橙框）
\newtcolorbox{pitfall}[1][常见陷阱]{
  enhanced, breakable, arc=3pt,
  before skip=10pt, after skip=10pt,
  colback=orange!6, colframe=orange!70!black,
  boxrule=0.7pt, left=8pt, right=8pt, top=6pt, bottom=6pt,
  fonttitle=\bfseries, title={【陷阱】#1},
}
% 追问 / Follow-up（蓝框）
\newtcolorbox{followup}[1][面试官追问]{
  enhanced, breakable, arc=3pt,
  before skip=10pt, after skip=10pt,
  colback=blue!4, colframe=blue!55!black,
  boxrule=0.7pt, left=8pt, right=8pt, top=6pt, bottom=6pt,
  fonttitle=\bfseries, title={【追问】#1},
}
% 复杂度 / 要点速记（绿框）
\newtcolorbox{keypoint}[1][要点]{
  enhanced, breakable, arc=3pt,
  before skip=10pt, after skip=10pt,
  colback=green!4, colframe=green!45!black,
  boxrule=0.7pt, left=8pt, right=8pt, top=6pt, bottom=6pt,
  fonttitle=\bfseries, title={【要点】#1},
}
% 面试话术：Code Review / 设计轮中"该怎么说出口"（紫框）
\newtcolorbox{sayit}[1][面试中这样说]{
  enhanced, breakable, arc=3pt,
  before skip=10pt, after skip=10pt,
  colback=violet!4, colframe=violet!60!black,
  boxrule=0.7pt, left=8pt, right=8pt, top=6pt, bottom=6pt,
  fonttitle=\bfseries, title={【话术】#1},
}

% ============================================================================
% 题目标题命令
% ============================================================================
\newcommand{\difficulty}[1]{%
  \IfStrEqCase{#1}{%
    {Easy}{\textcolor{duogreen!70!black}{\textbf{Easy}}}%
    {Medium}{\textcolor{duoorange!85!black}{\textbf{Medium}}}%
    {Hard}{\textcolor{duored!80!black}{\textbf{Hard}}}%
  }[\textbf{#1}]%
}

% 算法题：\problem{编号}{标题}{难度}
\newcommand{\problem}[3]{%
  \section{#2}%
  \noindent{\small\textcolor{abgray}{#1 \quad|\quad 难度：}\difficulty{#3}}\par
  \vspace{4pt}\hrule height 0.4pt \vspace{8pt}%
}

% 查错题：\bugcase{编号}{标题}{类别}
\newcommand{\bugcase}[3]{%
  \section{#2}%
  \noindent{\small\textcolor{abgray}{查错 \##1 \quad|\quad 类别：}\textcolor{duored!80!black}{\textbf{#3}}}\par
  \vspace{4pt}\hrule height 0.4pt \vspace{8pt}%
}

% 设计题：\designq{编号}{标题}{难度}
\newcommand{\designq}[3]{%
  \section{#2}%
  \noindent{\small\textcolor{abgray}{设计题 \##1 \quad|\quad 难度：}\difficulty{#3}}\par
  \vspace{4pt}\hrule height 0.4pt \vspace{8pt}%
}

% ============================================================================
% 信息可靠性标记：区分「官方」「社区面经」「高置信推断」
% ============================================================================
\newcommand{\srcofficial}{%
  \textcolor{duogreen!65!black}{\small\textbf{[官方]}}\ }
\newcommand{\srccommunity}{%
  \textcolor{duoorange!85!black}{\small\textbf{[面经报告]}}\ }
\newcommand{\srcinferred}{%
  \textcolor{duoblue!70!black}{\small\textbf{[高置信推断]}}\ }

% 常用记号
\newcommand{\bigO}[1]{\ensuremath{O(#1)}}
\newcommand{\code}[1]{\texttt{#1}}
